\documentclass[aspectratio=169]{beamer}
\usetheme{metropolis}
\metroset{sectionpage=none, numbering=counter, progressbar=frametitle}
\usepackage{amsmath, amssymb}
\usepackage[T1]{fontenc}

\title{Domain-Partition Algorithm for Turbomachinery Blades}
\subtitle{Slide 5: Frame Field}
\date{\today}
\author{}
\institute{}

\begin{document}

\begin{frame}{Frame Field --- Harmonic Laplace + Dirichlet}

\begin{center}
\Large
\[
\Delta \theta = 0 \quad \text{s.t.} \quad \theta|_{\partial\Omega} = \theta_{\text{wall}}
\]
\end{center}

\vspace{0.5em}
\begin{itemize}
    \item \textbf{Solve harmonic Laplace} for a smooth angle field inside the domain
    \item \textbf{Dirichlet BC:} boundary values align with wall cross-vectors (4-RoSy)
    \item \textbf{Normalize} to extract the 4-RoSy representative ($\mod \pi/2$)
\end{itemize}

\vspace{0.8em}
\small
\textit{Solver:} \texttt{tools/frame_field.py} \hfill
\textit{Ref:} Kn\"{o}ppel et al., "Globally optimal direction fields", SIGGRAPH 2013

\note{
Speaker notes:
The frame field is built in three steps. First, a harmonic Laplace solve gives the smoothest possible angle field in the interior. Second, Dirichlet boundary conditions lock the field to the wall cross-vectors computed from the surface geometry (see scripts/stage1/partition_surface.py:230--280). Third, normalization folds the continuous angle into the 4-RoSy representative space, i.e. modulo pi/2. The external solver in tools/frame_field.py performs the actual linear solve; this repo consumes the resulting field to detect singularities and integrate streamlines.
}

\end{frame}

\end{document}
